% GENERATED FILE. Edit canonical lessons, then run npm run generate:books.
% canonical-source-hash: fnv1a64:9e8fa0aebfa5840f
% canonical-lessons: HI-C551-manzur, HI-C551-azad, HI-C551-lacar, HI-C551-ariyal, HI-C551-zabardasti

\chapter{Free to Choose}
\label{ch:hi-free-to-choose}
\hlchaptermodality{551}

\begin{chapteropening}
\textbf{By the end of this chapter:} \emph{I can say whether someone is free, stubborn or shy.}

Complete the last lesson of chapter 551: I can say whether someone is free, stubborn or shy.
\end{chapteropening}

\section[\texorpdfstring{\hi{मंज़ूर} (manzūr)}{मंज़ूर (manzūr)}]{\hi{मंज़ूर} (manzūr) — approved; accepted}

\label{lesson:HI-C551-manzur}



\begin{quote}
Before the new one: say the Hindi for favourite, then the Hindi for a selection.
\end{quote}

\subsection*{You'll want to know: \hi{मंज़ूर}}
\textbf{\hi{मंज़ूर}} — \emph{manzūr} — \textquotedblleft{}approved; accepted\textquotedblright{}.

\begin{quote}
\hi{मालिक} \hi{को} \hi{मेरी} \hi{छुट्टी} \hi{मंज़ूर} \hi{है।}
\end{quote}

\begin{quote}
\emph{Mālik ko merī chuṭṭī manzūr hai.}
\end{quote}

\begin{quote}
The owner has approved my leave.
\end{quote}

\begin{cousinweb}
\textbf{\hi{मंज़ूर}} (\emph{manzūr}) is usually traced to Arabic \emph{manẓūr}, \textquotedblleft{}seen, looked upon\textquotedblright{}, through Persian. Persian and Hindi use it for something approved.
\end{cousinweb}

\begin{grammarlens}[title={What you've built}]
One more word for saying what happens if.
\end{grammarlens}

\subsection*{Your turn}
\begin{itemize}
\raggedright
  \item \emph{Say it:} \emph{manzūr}
  \item \emph{Say it:} \emph{manzūr}, once more
  \item \emph{Say it:} \emph{Mālik ko merī chuṭṭī manzūr hai.} — The owner has approved my leave.
  \item \emph{Recall:} say the Hindi for favourite, then the Hindi for a selection, then say \emph{manzūr} again
\end{itemize}

\begin{tcolorbox}[breakable,colback=teal!4,colframe=teal!35!black,title={Before you move on}]
What does \hi{मंज़ूर} mean? (\textquotedblleft{}Approved; accepted\textquotedblright{}.) Say it once more.
\end{tcolorbox}
\section[\texorpdfstring{\hi{आज़ाद} (āzād)}{आज़ाद (āzād)}]{\hi{आज़ाद} (āzād) — free}

\label{lesson:HI-C551-azad}



\begin{quote}
Before the new one: say the Hindi for a selection, then the Hindi for approved.
\end{quote}

\subsection*{You'll want to know: \hi{आज़ाद}}
\textbf{\hi{आज़ाद}} — \emph{āzād} — \textquotedblleft{}free\textquotedblright{}.

\begin{quote}
\hi{स्कूल} \hi{के} \hi{बाद} \hi{बच्चे} \hi{खेलने} \hi{के} \hi{लिए} \hi{आज़ाद} \hi{हैं।}
\end{quote}

\begin{quote}
\emph{Skūl ke bād bacce khelne ke liye āzād haiṁ.}
\end{quote}

\begin{quote}
After school the children are free to play.
\end{quote}

\begin{cousinweb}
\textbf{\hi{आज़ाद}} (\emph{āzād}) is usually traced to Persian \emph{āzād}, \textquotedblleft{}free\textquotedblright{}.
\end{cousinweb}

\begin{grammarlens}[title={What you've built}]
One more word for saying what happens if.
\end{grammarlens}

\subsection*{Your turn}
\begin{itemize}
\raggedright
  \item \emph{Say it:} \emph{āzād}
  \item \emph{Say it:} \emph{āzād}, once more
  \item \emph{Say it:} \emph{Skūl ke bād bacce khelne ke liye āzād haiṁ.} — After school the children are free to play.
  \item \emph{Recall:} say the Hindi for a selection, then the Hindi for approved, then say \emph{āzād} again
\end{itemize}

\begin{tcolorbox}[breakable,colback=teal!4,colframe=teal!35!black,title={Before you move on}]
What does \hi{आज़ाद} mean? (\textquotedblleft{}Free\textquotedblright{}.) Say it once more.
\end{tcolorbox}
\section[\texorpdfstring{\hi{लाचार} (lācār)}{लाचार (lācār)}]{\hi{लाचार} (lācār) — helpless}

\label{lesson:HI-C551-lacar}



\begin{quote}
Before the new one: say the Hindi for approved, then the Hindi for free.
\end{quote}

\subsection*{You'll want to know: \hi{लाचार}}
\textbf{\hi{लाचार}} — \emph{lācār} — \textquotedblleft{}helpless\textquotedblright{}.

\begin{quote}
\hi{पैर} \hi{टूटने} \hi{के} \hi{बाद} \hi{वह} \hi{कुछ} \hi{दिन} \hi{लाचार} \hi{था।}
\end{quote}

\begin{quote}
\emph{Pair ṭūṭne ke bād vah kuch din lācār thā.}
\end{quote}

\begin{quote}
After breaking his leg he was helpless for a few days.
\end{quote}

\begin{cousinweb}
\textbf{\hi{लाचार}} (\emph{lācār}) is usually traced to Persian \emph{lāchār}, \textquotedblleft{}without remedy, helpless\textquotedblright{}.
\end{cousinweb}

\begin{grammarlens}[title={What you've built}]
One more word for saying what happens if.
\end{grammarlens}

\subsection*{Your turn}
\begin{itemize}
\raggedright
  \item \emph{Say it:} \emph{lācār}
  \item \emph{Say it:} \emph{lācār}, once more
  \item \emph{Say it:} \emph{Pair ṭūṭne ke bād vah kuch din lācār thā.} — After breaking his leg he was helpless for a few days.
  \item \emph{Recall:} say the Hindi for approved, then the Hindi for free, then say \emph{lācār} again
\end{itemize}

\begin{tcolorbox}[breakable,colback=teal!4,colframe=teal!35!black,title={Before you move on}]
What does \hi{लाचार} mean? (\textquotedblleft{}Helpless\textquotedblright{}.) Say it once more.
\end{tcolorbox}
\section[\texorpdfstring{\hi{अड़ियल} (aṛiyal)}{अड़ियल (aṛiyal)}]{\hi{अड़ियल} (aṛiyal) — stubborn; obstinate}

\label{lesson:HI-C551-ariyal}



\begin{quote}
Before the new one: say the Hindi for free, then the Hindi for helpless.
\end{quote}

\subsection*{You'll want to know: \hi{अड़ियल}}
\textbf{\hi{अड़ियल}} — \emph{aṛiyal} — \textquotedblleft{}stubborn; obstinate\textquotedblright{}.

\begin{quote}
\hi{मेरा} \hi{गधा} \hi{बहुत} \hi{अड़ियल} \hi{है।}
\end{quote}

\begin{quote}
\emph{Merā gadhā bahut aṛiyal hai.}
\end{quote}

\begin{quote}
My donkey is very stubborn.
\end{quote}

\begin{cousinweb}
\textbf{\hi{अड़ियल}} (\emph{aṛiyal}) is built on \emph{aṛnā}, \textquotedblleft{}to come to a stop, to dig one's heels in\textquotedblright{}, which CDIAL files under \emph{*aḍ}, \textquotedblleft{}obstruct, stop\textquotedblright{} (CDIAL 187), with the ending -iyal for someone given to it.
\end{cousinweb}

\begin{grammarlens}[title={What you've built}]
One more word for saying what happens if.
\end{grammarlens}

\subsection*{Your turn}
\begin{itemize}
\raggedright
  \item \emph{Say it:} \emph{aṛiyal}
  \item \emph{Say it:} \emph{aṛiyal}, once more
  \item \emph{Say it:} \emph{Merā gadhā bahut aṛiyal hai.} — My donkey is very stubborn.
  \item \emph{Recall:} say the Hindi for free, then the Hindi for helpless, then say \emph{aṛiyal} again
\end{itemize}

\begin{tcolorbox}[breakable,colback=teal!4,colframe=teal!35!black,title={Before you move on}]
What does \hi{अड़ियल} mean? (\textquotedblleft{}Stubborn; obstinate\textquotedblright{}.) Say it once more.
\end{tcolorbox}
\section[\texorpdfstring{\hi{ज़बरदस्ती} (zabardastī)}{ज़बरदस्ती (zabardastī)}]{\hi{ज़बरदस्ती} (zabardastī) — force, coercion}

\label{lesson:HI-C551-zabardasti}



\begin{quote}
Before the new one: say the Hindi for helpless, then the Hindi for stubborn.
\end{quote}

\subsection*{You'll want to know: \hi{ज़बरदस्ती}}
\textbf{\hi{ज़बरदस्ती}} — \emph{zabardastī} — \textquotedblleft{}force, coercion\textquotedblright{}.

\begin{quote}
\hi{बच्चे} \hi{से} \hi{ज़बरदस्ती} \hi{मत} \hi{कीजिए।}
\end{quote}

\begin{quote}
\emph{Bacce se zabardastī mat kījiye.}
\end{quote}

\begin{quote}
Don't force the child.
\end{quote}

Feminine.

\begin{cousinweb}
\textbf{\hi{ज़बरदस्ती}} (\emph{zabardastī}) is usually traced to Persian \emph{zabardastī}, \textquotedblleft{}force, having the upper hand\textquotedblright{}.
\end{cousinweb}

\begin{grammarlens}[title={What you've built}]
One more word for saying what happens if.
\end{grammarlens}

\subsection*{Your turn}
\begin{itemize}
\raggedright
  \item \emph{Say it:} \emph{zabardastī}
  \item \emph{Say it:} \emph{zabardastī}, once more
  \item \emph{Say it:} \emph{Bacce se zabardastī mat kījiye.} — Don't force the child.
  \item \emph{Recall:} say the Hindi for helpless, then the Hindi for stubborn, then say \emph{zabardastī} again
\end{itemize}

\begin{tcolorbox}[breakable,colback=teal!4,colframe=teal!35!black,title={Before you move on}]
What does \hi{ज़बरदस्ती} mean? (\textquotedblleft{}Force, coercion\textquotedblright{}.) Say it once more.
\end{tcolorbox}
